\chapter{Runtime architecture}

The runtime data flow has three tiers. Tier one is the player's browser:
the federated Angular host, the active remote module, and the WASM iframe.
Tier two is the per-remote nginx fleet: the host, four route-or-game
remotes, and the asset-tier \texttt{wasm-fragment}. Tier three is the BFF
plus its persistence: Spring Boot WebFlux on top of Postgres and Cassandra,
behind Casdoor for authentication.

\section{Browser-side flow}

\begin{figure}[!h]
\centering
\snapdiagram{diagrams/runtime-flow.tex}
\caption{Runtime data flow from a click on the home page to a frame
rendered inside the Godot iframe. \texttt{SCENE\_TYPE\_TO\_GAME\_SLUG}
forks the renderer; federated games mint a \texttt{launch\_uid} for
operational visibility; \texttt{TW\_LEVEL\_TRANSITION} \texttt{postMessage}
is caught at the host level so it works regardless of which renderer is
mounted.}
\end{figure}

The host shell's \snapfile{main.ts} is intentionally split into two
modules. \snapfile{main.ts} runs the
\textbf{native-federation \texttt{initFederation(...)}} bootstrap with the
manifest fetched from \texttt{GET /api/federation/manifest}, then stashes
\texttt{loadRemoteModule} on \texttt{globalThis.\_\_SNAPDRAGON\_LOAD\_REMOTE\_MODULE\_\_}
and only then dynamically imports \snapfile{bootstrap.ts}, which calls
\texttt{bootstrapApplication} with the host's \texttt{AppComponent}. This
ordering is non-negotiable: any \texttt{@angular/*} static import resolved
before \texttt{initFederation} runs will fail because the import map
that's about to be installed isn't there yet.

\begin{snapnote}[title={Why \texttt{globalThis} and not a service?}]
The federation runtime can't be reached through Angular's DI without
making \texttt{@angular-architects/native-federation} a static import
into a route-handler chunk, which puts it on karma's module graph and
breaks tests that don't load the federation runtime at all. Stashing the
function on \texttt{globalThis} after bootstrap keeps the package out of
karma's static graph and out of every route bundle that doesn't need it.
\end{snapnote}

\section{Three-tier topology}

\begin{figure}[!h]
\centering
\snapdiagram{diagrams/topology.tex}
\caption{Compose-network view. Each box is a Docker container; arrows are
HTTP routes proxied by the host nginx at \snapfile{frontend/nginx.conf}.
The BFF is the only container that talks to Postgres / Cassandra /
Casdoor.}
\end{figure}

\subsection{The host shell}

The host nginx serves the host bundle plus a set of \texttt{location
/\_remotes/<slug>/} blocks that proxy each remote's
\texttt{remoteEntry.json} and chunks. For \texttt{wasm-fragment} the host
strips upstream COOP/COEP duplicates with \texttt{proxy\_hide\_header} and
re-emits the canonical set so the WASM iframe loads with isolated
headers.

\subsection{The remotes}

\begin{tabularx}{\linewidth}{l l X}
\toprule
\textbf{Slug} & \textbf{Kind} & \textbf{What it exposes / does} \\
\midrule
\texttt{login} & \texttt{remote-shell} & \texttt{./LoginComponent} —
mounted via \texttt{loadComponent} on the \texttt{/login} route. Inline
\texttt{POST /api/auth/token} + \texttt{sessionStorage.tw\_auth} write,
mirrors the host \texttt{AuthService.StoredAuth} shape. \\
\texttt{admin} & \texttt{remote-shell} & \texttt{./AdminShell} +
\texttt{./UploadShell} — both expose under
\texttt{federation\_modules.slug = 'admin'} so \texttt{/admin} and
\texttt{/admin/games} share a build. \\
\texttt{games-three-wizards} & \texttt{remote-game} &
\texttt{./GameRoot} — Web Component (\texttt{createCustomElement}) with
\texttt{@Input levelSlug + levelId + launchUid}. Renders cave and field. \\
\texttt{games-dodecahedron} & \texttt{remote-game} &
\texttt{./GameRoot} — same pattern, single chrome-less scene. \\
\texttt{wasm-fragment} & \texttt{fragment} & Asset-tier nginx; bakes
\snapfile{backend/src/main/resources/assets/levels/} and serves at
\texttt{/\_remotes/wasm-fragment/<slug>/index.html}. \\
\bottomrule
\end{tabularx}

\subsection{The BFF}

\snapfile{backend/} runs as a single Spring Boot 3.4 application with both
WebFlux (R2DBC, runtime path) and JDBC (Liquibase, blocking) wired in.
Spring auto-config skips \texttt{DataSourceAutoConfiguration} when an
R2DBC \texttt{ConnectionFactory} is present, so
\snapfile{JdbcDataSourceConfig} declares an explicit
\texttt{@Primary HikariDataSource} of pool size 3 just for Liquibase to
acquire on boot.

\section{Authentication and access}

Casdoor is reachable only from inside the compose network. The host BFF
proxies the login flow at \texttt{/api/auth/*} and verifies JWTs against
the anonymous \texttt{/.well-known/jwks} endpoint (not
\texttt{/api/get-jwks}, which requires admin auth and would 401 on every
request). Per-game visibility is enforced by the \texttt{game\_access}
table — see chapter~\ref{ch:datamodel}.

\section{Server-sent events}

\snapfile{GameStreamController} merges \emph{eight} SSE channels onto a
single \texttt{Flux<ServerSentEvent>}: seven \texttt{Sinks.Many} from
\snapfile{GameStateService} (\texttt{movement}, \texttt{world},
\texttt{character}, \texttt{scene}, \texttt{input}, \texttt{audio},
\texttt{wasm-console}) plus one from \snapfile{GameTimeService}
(\texttt{time}), with a 15-second \texttt{Flux.interval} keepalive
that fires immediately on subscribe so the response headers flush at
connect time and the browser \texttt{EventSource.onopen} fires in
milliseconds.

The latter three channels are part of the generalised input surface
introduced for gamepad / mouse / keyboard / touch / audio support
(see chapter~\ref{ch:datamodel}'s reference-table description and
chapter~\ref{ch:api}'s \emph{Input events} + \emph{Diagnostics} tags).
Subscribers opt into individual channels via the EventSource
event-listener model (\texttt{source.addEventListener('input', cb)});
merging onto the single connection keeps the host page from needing
N EventSources to see everything.

\begin{snapcaution}[title={The \texttt{autoCancel} trap}]
\texttt{Sinks.many().multicast().onBackpressureBuffer(N)} defaults to
\texttt{autoCancel=true}. When every subscriber disconnects (iframe
remount, EventSource reconnect, network blip), the sink terminates
\textbf{permanently} and all subsequent \texttt{tryEmitNext} calls
silently drop. Every sink in this project passes \texttt{autoCancel=false}
explicitly. If a freshly-deployed sink \enquote{works once and then
silently dies}, this is the cause.
\end{snapcaution}

\section{Per-session SSE sink isolation}

Multiplayer sessions add a per-session sink layer on top of the
legacy global sinks. \snapfile{GameStateService} maintains a
\texttt{ConcurrentHashMap<String, SessionSinks>} lazy-allocated by
\texttt{sessionId}; when a control event arrives with a session id,
\texttt{applyControl(event, characterId, sessionId)} fans out to the
matching per-session bundle \emph{and} the global sink. A subscriber
to \texttt{GET /api/game/stream?sessionId=$\alpha$} sees only session
$\alpha$'s events; the admin firehose at
\texttt{/api/admin/stream} keeps the cross-session view it had before.

\begin{figure}[!h]
\centering
\snapdiagram{diagrams/sessions-topology.tex}
\caption{Per-session SSE sink isolation. Events emitted on session
$\alpha$'s sinks NEVER reach session $\beta$'s subscribers and vice
versa; the global sink keeps the admin firehose intact. The
\reposrc{backend/src/test/java/games/snapdragon/bff/service/SessionIsolationTest.java}
integration test pins the dashed-arrow non-leakage in both directions.}
\end{figure}

The new \emph{view} channel carries per-player camera + orientation
samples ($\sim$10\,Hz) emitted by the Rust GDExtension's
\texttt{view\_egress} signal. Admins watch via
\texttt{GET /api/admin/sessions/\{id\}/stream} which merges movement,
world, character, and view onto a single per-session connection.
Without per-session keying, two concurrent sessions on the same pod
would cross-contaminate every fan-out — so the keyed map is the
load-bearing piece.

\section{Maker side services}

The RPG and VN makers each ship as their own reactive Spring Boot
service alongside the BFF, with their own Postgres tenants:

\begin{itemize}
  \item \snapfile{rpg-builder/} --- Mudlet-shaped editor data model
        (projects, areas, rooms, exits, items, NPCs, scripts) backed
        by the \texttt{postgres-rpg-builder} container. REST surface
        at \texttt{/api/rpg/*}.
  \item \snapfile{vn-builder/} --- godot-2d-visual-novel-shaped data
        model (projects, acts, characters, backgrounds, scenes, lines,
        choices) backed by the \texttt{postgres-vn-builder} container.
        REST surface at \texttt{/api/vn/*}.
\end{itemize}

Each side service's build worker walks the project tree, emits a Godot
4 source tree with \snapfile{snapdragon\_bridge.gdextension} baked in
(multiplayer-aware out of the box), runs the Godot CLI per area or act,
zips the result, and POSTs to the BFF's
\texttt{/api/admin/games/upload} endpoint \emph{with the originating
publisher's bearer token}. The BFF's existing
\texttt{publisher\_game\_scopes} check enforces who can ship which
slug, so a single compromised maker can't spam the catalog. Decision
018 covers the side-service architecture; the per-database schemas
live in the \emph{Schema diagrams} appendix at the end of this
document.

\section{Mod / publisher per-game scope}

Authorisation for mod- and publisher-only endpoints reads two layers,
both required:

\begin{itemize}
  \item Casdoor membership in \texttt{ROLE\_MOD} or
        \texttt{ROLE\_PUBLISHER} (claim is in the JWT).
  \item One row in \texttt{mod\_game\_scopes} /
        \texttt{publisher\_game\_scopes} per (casdoorId, gameSlug)
        pair (Postgres).
\end{itemize}

Admins manage both layers via \texttt{/api/admin/role-grants/*};
publishers can grant mod on \emph{games they own} via
\texttt{/api/publisher/grant-mod}, which walks the same write path but
enforces caller-owns-target-game before mutating. The full rejection
matrix is pinned by
\reposrc{e2e/features/security/mod-grant-gates.feature} on the BDD side
and a \texttt{@WebFluxTest} on the controller side.
